% Seção: Avaliação de desempenho
% Dono: Rafael · Orçamento: 1 página · docs/fontes/plano-de-projeto.md §14.1
% Conteúdo: complexidade teórica e tempos medidos de cada algoritmo; comparação com
% NetworkX só se a etapa opcional (F-AVAL-OPC) acontecer.
% A tabela é preenchida na F-AVAL (RF-AVA-02) com os números registrados nos arquivos
% de feature. Nenhum número entra aqui sem o comando que o gerou.
\section{Avaliação de desempenho}\label{sec:avaliacao}

\subsection{Metodologia de medição}

Todos os tempos deste relatório foram obtidos com a mesma função,
\texttt{medir}, do módulo \texttt{transporte.util.medicao}. Ela executa o
algoritmo $N$ vezes com os mesmos argumentos (por padrão, $N = 5$) e
cronometra cada execução isoladamente com \texttt{time.perf\_counter}, o
relógio monotônico de maior resolução do Python~\cite{python-perf-counter}.
Só a chamada do algoritmo fica dentro do intervalo medido: a leitura do GTFS e
a construção do grafo são medidas à parte, quando interessam. O resultado é a
média $\bar{t}$ e o desvio padrão amostral
\[
  s = \sqrt{\frac{1}{N-1}\sum_{i=1}^{N}(t_i - \bar{t})^2},
\]
com divisor $N - 1$ porque as $N$ execuções são uma amostra do tempo do
algoritmo naquela máquina, e não a população inteira. Com $N = 1$, o desvio é
reportado como zero.

Cada tempo é acompanhado do tamanho do grafo ($|V|$ e $|E|$ da visão usada) e
da máquina em que foi medido, porque tempo de parede só é comparável dentro da
mesma máquina e da mesma rede. A complexidade teórica de cada algoritmo vem da
docstring da função correspondente, justificada no arquivo da feature.

\subsection{Complexidade e tempos medidos}

\pendente{Tabela consolidada (F-AVAL, RF-AVA-02): algoritmo, visão, complexidade
teórica, $|V|$, $|E|$, tempo médio $\pm$ desvio e $N$, com os números
registrados por cada dono na rede real.}

\subsection{Comparação com o NetworkX}

\pendente{Só se a etapa opcional F-AVAL-OPC acontecer (16/11 a 18/11); caso
contrário, esta subseção sai.}
